\chapter*{前言：这本书不教你写代码}
\addcontentsline{toc}{chapter}{前言：这本书不教你写代码}

本书围绕一条主线展开：

\begin{quote}
AI 负责生成候选答案，人负责定义正确答案的标准。
\end{quote}

\section*{为什么现在需要一种新的工作方式}

设想一名交通工程专业的学生，手中有一份出租车 GPS 数据、一个交通流预测任务，以及一个并不熟悉的 Python 环境。按照传统方式，他可能需要花几天时间查找工具、编写代码并处理报错；在 AI 协作下，他可以更快完成数据读取、清洗、建模、可视化和初步报告。

但速度不等于正确。AI 给出的模型可能拥有很高的决定系数，却因为随机划分时序数据而造成严重的数据泄露。AI 能让人更快到达一个答案，但不能保证那个答案是对的。

\section*{Vibe Coding到底是什么}

Vibe Coding 并不意味着不需要理解代码，而是不再需要从零开始实现每一行代码。AI 无法自动判断研究问题是否合理，不知道交通数据中的异常值意味着什么，也无法独立识别所有逻辑漏洞和数据泄露风险。

人的作用正在从代码的逐行实现，转向目标定义、约束控制和结果判断。

\section*{本书的主线}

\begin{enumerate}
  \item 如何提问（Prompting）。
  \item 如何建模（Modeling）。
  \item 如何判断（Evaluation）。
\end{enumerate}

\section*{如何使用本书}

\begin{itemize}
  \item 本科生可以将本书作为使用 AI 辅助完成课程项目的指南。
  \item 研究生可以将本书作为使用 AI 辅助科研建模与实验设计的手册。
  \item 工程师可以将本书作为快速完成业务分析原型的工具书。
\end{itemize}
